\section{Operator Programs}
\label{sec:operator-programs}

An operator of arity (k) acts within the space of valid Abstract Graphs,
\begin{equation}
  o:\mathcal{A}^{k}\rightharpoonup\mathcal{A},
\end{equation}
subject to its declared input contracts. Unary and multi-input operators can
therefore form expressions without leaving the Abstract Graph representation.
For example, a sequential unary program is
\begin{equation}
  P=o_k\circ\cdots\circ o_1.
\end{equation}

This closure is a well-formedness and expressivity property, not a refinement
order: applying or composing operators modifies the representation and need
not monotonically increase discrimination.

This section will define operator contracts and typing, composition rules,
representation invariants, the reference operator suite, and computational
complexity. The principal invariants are mapping validity, provenance
completeness, permutation invariance, determinism, composition closure,
serialization round-trip preservation, and identity consistency.

The operational extension point is deliberately small. A domain expert may
provide a Python function from a NetworkX graph to a concrete list of node
groups. A constructor materializes those groups and their mappings as the
output Abstract Graph. Disjoint exhaustive groups define a partition, while
operators may declare overlapping or partial covers. The same contract also
makes generated ad hoc operators possible, provided they pass validation for
determinism, invariance, coverage, semantics, and resource use.
